\section{Weighted polynomial norms: an arithmetic zero-count engine}
\label{sec:weighted-norm}

The previous framework packet proves a sharp zero count for equations
$Y^2=f(X)$ with $\deg f$ odd \cite{repoEmbedding}. The construction below extends that
mechanism to an entire class of plane curves, including equations with
mixed $X,Y$ terms.  It also gives an exact dimension formula for a family
of finite spaces of functions.  The arguments require only polynomial
arithmetic, determinants, exact rational comparisons, and induction on
bounded lists.  They use neither a divisor theorem nor a choice of
algebraic branches.

\subsection{Input data and the exact norm degree}

Fix coprime positive integers $m,d$ with $m\geq2$, and a nonzero rational
$c$.  Give $X$ weight $m$ and $Y$ weight $d$.  Consider
\begin{equation}\label{eq:weighted-curve}
 F(X,Y)=Y^m-cX^d+
       \sum_{mi+dj<md} f_{ij}X^iY^j,
 \qquad f_{ij}\in\mathbb Q.
\end{equation}
The sum is a finite coefficient table; missing coefficients mean zero.
All its $Y$ exponents are less than $m$.  Consequently $F$ is monic of
$Y$ degree $m$ and the quotient
\[
 \mathcal A=\mathbb Q[X,Y]/(F)
\]
is a free $\mathbb Q[X]$ module with the explicit basis
$1,Y,\ldots,Y^{m-1}$.  Here freeness is proved directly by division by a
monic polynomial: subtract the leading multiple of $F$ until the $Y$
degree is below $m$; a nonzero multiple of the monic $F$ cannot itself
have $Y$ degree below $m$.

A reduced polynomial is an expression
\[
 G(X,Y)=\sum_{j<m} A_j(X)Y^j.
\]
Suppose $G\ne0$ and let
\begin{equation}\label{eq:weighted-size}
 W(G)=\max\{mi+dj:[X^iY^j]G\ne0\}.
\end{equation}
Write $M_G(X)$ for the $m$ by $m$ matrix of multiplication by $G$ on
$\mathcal A$ in the displayed basis, and put
\[
 N_G(X)=\det M_G(X).
\]
These are explicit polynomials.  To compute column $j$, reduce $GY^j$
by the equation $Y^m=cX^d-\sum f_{ij}X^iY^j$.  Its starting $Y$ degree
is at most $2m-2$; each reduction lowers the largest remaining $Y$
exponent.  The determinant can be defined by its finite permutation sum.
This definition is useful even when a faster elimination algorithm is
used by an implementation.

\begin{theorem}[Exact weighted norm degree]\label{thm:weighted-norm-degree}
There is a unique monomial $aX^kY^j$ of $G$ of maximum weight, with
$a\ne0$ and $0\le j<m$.  Its norm satisfies
\[
 \deg N_G=W(G)=mk+dj,
 \qquad
 \operatorname{lc}(N_G)
   =a^m\bigl((-1)^{m-1}c\bigr)^j\ne0.
\]
In particular the matrix determinant is nonzero, including when the curve
has singular affine points.
\end{theorem}

\begin{proof}
If $mi+dj=mi'+dj'$ with $0\le j,j'<m$, then $m$ divides
$d(j-j')$.  Coprimality implies that $m$ divides $j-j'$, and its absolute
value is less than $m$.  Thus $j=j'$ and $i=i'$.  Distinct reduced
monomials therefore have distinct weights.  This proves the uniqueness
of the maximum.

The leading determinant can be calculated without guessing which of its
summands cancel.  Introduce an indeterminate $t$ and define
\begin{align*}
 F_t(Z)&=t^{md}F(t^{-m},t^{-d}Z),\\
 G_t(Z)&=t^{W(G)}G(t^{-m},t^{-d}Z).
\end{align*}
The weight inequalities show that both expressions lie in
$\mathbb Q[t,Z]$: a monomial of weight $v$ is multiplied by
$t^{md-v}$ or by $t^{W(G)-v}$.  Moreover $F_t$ is monic of degree $m$
in $Z$, and specialization at $t=0$ gives
\[
 F_0(Z)=Z^m-c,\qquad G_0(Z)=aZ^j.
\]
Let $D(t)$ be the determinant of multiplication by $G_t$ on
$\mathbb Q[t,Z]/(F_t)$ in its monic-remainder basis.  This is a polynomial
in $t$.  Over the Laurent polynomial ring, the change of basis
$Y=t^{-d}Z$ is invertible, so conjugating the multiplication matrices and
factoring the scalar $t^{-W(G)}$ out of their $m$ columns gives
\begin{equation}\label{eq:weighted-rescaling}
 D(t)=t^{mW(G)}N_G(t^{-m}).
\end{equation}
This identity can alternatively be verified directly on the $m$ basis
vectors; no semantic assertion about a field of algebraic functions is
needed.

At $t=0$, multiplication by $Z$ in
$\mathbb Q[Z]/(Z^m-c)$ has the cyclic companion matrix whose determinant
is $(-1)^{m-1}c$.  Multiplication by $aZ^j$ is $a$ times the $j$th power
of that matrix.  Determinant multiplicativity therefore yields
\[
 D(0)=a^m\bigl((-1)^{m-1}c\bigr)^j\ne0.
\]
Write $N_G(X)=\sum_{v=0}^{q}n_vX^v$ with $n_q\ne0$, if the
polynomial is nonzero.  Equation~\eqref{eq:weighted-rescaling} first
shows that it cannot be zero.  Its right side has terms
$n_vt^{m(W(G)-v)}$.  Distinct $v$ have distinct exponents.  Since the
left side has no negative powers, $q\le W(G)$.  Since its constant
coefficient is nonzero, $q=W(G)$, and that coefficient is the displayed
leading coefficient of $N_G$.  This proves all assertions.
\end{proof}

\paragraph{An immediate algebraic benefit.}
The quotient $\mathcal A$ has no zero divisors.  Indeed, for a nonzero
reduced $G$, the adjugate identity says
$\operatorname{adj}(M_G)M_G=N_G I$.  If $GH=0$, it follows that every
coefficient of the reduced $H$ is annihilated by the nonzero polynomial
$N_G$.  Rational polynomial arithmetic has no zero divisors, so $H=0$.
Thus an exact function-field interface can be built from pairs of
reduced polynomials with a nonzero denominator.  Equality is tested by
reducing cross products, and operations are finite polynomial
computations.  This observation supplies an algebraic interface; it does
not claim a general implementation of divisors or Riemann--Roch spaces.

\subsection{Counting points with repeated first coordinates}

A determinant root bound by itself counts distinct $X$ coordinates, not
points.  The following argument records the necessary multiplicities.

\begin{lemma}[A fiber contributes its full cardinality]\label{lem:weighted-fiber}
Fix $u\in\mathbb Q$.  If the distinct rationals
$v_1,\ldots,v_r$ satisfy
\[
 F(u,v_i)=G(u,v_i)=0\qquad(1\le i\le r),
\]
then $(X-u)^r$ divides $N_G(X)$.
\end{lemma}

\begin{proof}
The polynomial $F(u,Y)$ is monic of degree $m$, so the elementary
polynomial root bound gives $r\le m$.  Define row vectors
\[
 e(v_i)=(1,v_i,\ldots,v_i^{m-1}).
\]
Evaluation is compatible with reduction modulo $F(u,Y)$, because
$F(u,v_i)=0$.  Applying the evaluation row to each column of the
multiplication matrix therefore gives
\[
 e(v_i)M_G(u)=G(u,v_i)e(v_i)=0.
\]
Form an $m$ by $m$ rational matrix $U$ whose first $r$ rows are these
vectors and whose remaining rows are the standard basis rows with
indices $r,r+1,\ldots,m-1$.  Its determinant is the determinant of the
first $r$ columns of its first $r$ rows, namely
\[
 \det U=\prod_{1\le i<j\le r}(v_j-v_i)\ne0.
\]
The empty product is one.  Thus no unspecified basis extension is used.

The first $r$ rows of $UM_G(X)$ vanish at $X=u$.  The factor theorem
makes every entry of those rows divisible by $X-u$.  Factoring this
factor from each of the $r$ rows shows that $(X-u)^r$ divides
$\det(UM_G)=(\det U)N_G$.  Since $\det U$ is a nonzero rational,
it is invertible, and the same divisibility holds for $N_G$.
\end{proof}

\begin{theorem}[Sharp arithmetic zero count]\label{thm:weighted-zero-count}
Every finite list of pairwise distinct rational points satisfying
$F(x,y)=G(x,y)=0$ has length at most $W(G)$.
\end{theorem}

\begin{proof}
Sort the given finite list by its rational first coordinate, and let
$u_1,\ldots,u_s$ be the distinct first coordinates, with fiber sizes
$r_1,\ldots,r_s$.  Inside each fiber the second coordinates are distinct,
because the points are distinct.  Lemma~\ref{lem:weighted-fiber} gives
$(X-u_i)^{r_i}\mid N_G$ for every $i$.

The powers belonging to different $u_i$ are coprime.  For completeness,
this can be checked without factoring a general polynomial.  If
$P=(X-u)^r Q$ and $(X-v)^s\mid P$ with $u\ne v$, write the identity
in powers of $T=X-v$.  The first factor has the nonzero constant
coefficient $(v-u)^r$.  If the first $s$ coefficients of $P(v+T)$
vanish, the constant coefficient equation first forces $Q(v)=0$;
the coefficient equation at $T^k$ then forces the $k$th coefficient
of $Q(v+T)$ to vanish once the earlier ones vanish.  Induction on
$k<s$ proves $(X-v)^s\mid Q$.  Repeating this removal for the finite
list of fibers proves
\[
 \prod_{i=1}^s(X-u_i)^{r_i}\mid N_G.
\]
The product has degree $\sum_i r_i$, the length of the input list.
Theorem~\ref{thm:weighted-norm-degree} says that the nonzero polynomial
$N_G$ has degree $W(G)$, so $\sum_i r_i\le W(G)$.
\end{proof}

The proof handles ramification and singular fibers automatically.  It does
not assume that $F(u,Y)$ has $m$ distinct roots, that it splits over the
rationals, or that the roots can be selected continuously.

\subsection{Finite spaces and an elementary dimension formula}

For a numerical budget $L\ge0$ define the explicit monomial list
\[
 \mathcal B_L=\{X^iY^j:0\le j<m,\ i\ge0,\ mi+dj\le L\}.
\]
Every entry has $i\le L$ and $j<m$, so a rectangular bounded scan
constructs it.  Its elements are linearly independent in $\mathcal A$
by the monic-remainder uniqueness proved above.  Hence the dimension of
their span equals the directly computable integer
\begin{equation}\label{eq:weighted-dimension}
 n(L)=\sum_{j=0}^{m-1}
       \max\left(0,1+\left\lfloor\frac{L-dj}{m}\right\rfloor\right).
\end{equation}
In this displayed formula the floor is the signed integer floor.  A
natural-number program uses the equivalent branch: contribute zero if
$L<dj$, and otherwise contribute $1+\lfloor(L-dj)/m\rfloor$.

\begin{proposition}[Exact eventual dimension]\label{prop:weighted-dimension}
Put $g_0=(m-1)(d-1)/2$.  This is an integer.  For every
$L\ge\max(0,md-m-d)$,
\[
 n(L)=L+1-g_0.
\]
\end{proposition}

\begin{proof}
Since $m,d$ are coprime, they cannot both be even, so
$(m-1)(d-1)$ is even.  Under the stated lower bound on $L$,
$L-dj\ge -m$ for all $j<m$.  Therefore each summand
$1+\lfloor(L-dj)/m\rfloor$ is nonnegative and the maximum with zero
can be removed.

The residues $r_j$ of $L-dj$ modulo $m$ are a permutation of
$0,1,\ldots,m-1$, by coprimality.  Write
$L-dj=mq_j+r_j$, $0\le r_j<m$.  Summing these $m$ equations gives
\[
 m\sum_{j<m}q_j
   =mL-\frac{dm(m-1)}2-\frac{m(m-1)}2.
\]
Adding $m$ to $\sum_jq_j$ proves the claimed formula.
Every calculation is an identity of finite sums; no geometric dimension
or genus theorem is invoked.
\end{proof}

The number $g_0$ agrees with familiar geometric formulas in suitable
smooth cases, but no identification with the genus is needed here.  The
proposition proves exactly the numerical dimension assertion stated.

For a rational point $P=(x,y)$ of the curve, let
$\Phi_L(P)$ be the vector of the values of the monomials in
$\mathcal B_L$, in the scan order.  Because $1\in\mathcal B_L$,
this vector is nonzero.

\begin{corollary}[Uniform bound over a proper subspace]\label{cor:weighted-subspace}
Let $V$ be a proper rational linear subspace of
$\mathbb Q^{n(L)}$, specified by a finite rational basis or matrix.
Any finite list of distinct curve points whose evaluation vectors lie
in $V$ has length at most $L$.
\end{corollary}

\begin{proof}
Exact Gaussian elimination produces a nonzero rational row vector
$\lambda$ annihilating $V$: reduce the finite system of linear
equations for the annihilator, choose a free variable, set it to one,
and solve backwards with witnessed nonzero pivots.  Properness means
that a free variable exists.  Define
\[
 G_\lambda(X,Y)=\sum_{b\in\mathcal B_L}\lambda_b b(X,Y).
\]
It is a nonzero reduced polynomial, because the monomials are distinct
and $\lambda\ne0$, and it has weight at most $L$.  All the given points
are zeros of $G_\lambda$.  Apply
Theorem~\ref{thm:weighted-zero-count}.
\end{proof}

\begin{corollary}[The remaining approximation contract is isolated]\label{cor:weighted-cover}
Let $S$ be a decidable class of points on the curve.  Suppose every
finite list of distinct points in $S$ has its $\Phi_L$ vectors covered
by at most $M$ supplied proper rational subspaces.  Then every such list
has length at most $ML$.
\end{corollary}

\begin{proof}
For each point scan the finite list of subspaces and assign it to the
first one containing its vector.  Subspace membership is decidable by
finite exact elimination, and the assumed cover proves that the scan
succeeds.  The points assigned to one subspace form a distinct sublist,
so the preceding corollary bounds that sublist by $L$.  The original
length is the sum of at most $M$ such lengths, hence at most $ML$.
\end{proof}

This resolves the zero-count side of a finite-cover argument for every
curve of the displayed weighted form.  Producing or proving the stated
approximation cover is a separate dependency.  Neither the norm theorem
nor the dimension formula produces that cover.  The conclusion is a
cardinality bound on arbitrary supplied lists, not a height bound or a
complete enumeration of all points.

\subsection{Examples and sharpness}

Take $m=3$, $d=4$ and
\[
 F(X,Y)=Y^3-X(X-1)(X-2)(X-3).
\]
All terms other than $Y^3-X^4$ have weight below $12$.  For $G=Y$,
$W(G)=4$ and
\[
 N_G(X)=X(X-1)(X-2)(X-3).
\]
The four points $(0,0),(1,0),(2,0),(3,0)$ are distinct common zeros.
Thus the zero-count bound is attained.  Here
$\mathcal B_4=(1,X,Y)$ in a suitable order, and the proper subspace
with last coordinate zero has exactly four points in its inverse image.

The same method includes mixed equations such as
\[
 Y^3=X^4+XY+X+1.
\]
The mixed monomial $XY$ has weight $7<12$.  Every nonzero linear form
$\alpha+\beta X+\gamma Y$ has at most four rational zeros on this
curve, and every nonzero combination of
$1,X,Y,X^2,XY,Y^2$ has at most eight.  These assertions follow from the
weight bounds, with no test of smoothness and no root isolation.
For $m=3,d=4$ the dimension formula gives $n(L)=L-2$ for $L\ge5$;
for example $\mathcal B_8$ has six elements.

\subsection{A broader finite certificate when coprimality fails}

The rescaling argument has a useful form that does not assume
$\gcd(m,d)=1$.  Allow any monic $F$ of $Y$ degree $m$ whose monomials
have weight at most $md$, and any reduced nonzero $G$ of maximum
weight $W$.  Form $F_t,G_t$ by the same formulas, and compute
\[
 C(F,G)=\det\left(\text{multiplication by }G_0
                 \text{ on }\mathbb Q[Z]/(F_0)\right).
\]
Here $F_0$ is still monic of degree $m$.  If the exact rational test
$C(F,G)\ne0$ succeeds, the identical proof gives
\[
 \deg N_G=W,\qquad
 \#\{\text{distinct supplied common zeros}\}\le W.
\]
Thus coprimality is a simple sufficient condition making this test
succeed for every nonzero reduced $G$ in our principal family.  For
more general weighted curves the finite initial determinant is a usable
certificate on individual linear forms.  Failure of the test means only
that this leading-degree certificate did not apply; it does not decide
the number of zeros or the irreducibility of the original curve.

\subsection{The iPRA translation and a Lean development path}

The theorem is uniform in finite rational coefficient tables and in the
natural numbers $m,d,L$.  It must not be represented by a hypothetical
universal evaluator for arbitrary primitive recursive programs.  Instead
fix the finitely many algorithms below and represent each by its own
primitive recursive definition.

\begin{enumerate}
\item Encode signed integers, normalized fractions, coefficient lists,
matrices and point lists.  Implement rational comparison and equality by
cross multiplication with positive denominators.  Prove the totalized
operations correct on the decidable valid-code predicate.
\item Define rational polynomial addition, multiplication, coefficient
lookup, degree, evaluation, monic division in $Y$, and the finite
permutation determinant.  Prove their coefficient identities with
bounded folds.  A determinant specialization lemma and determinant
multiplicativity are explicit finite identities, not imported reflection
principles.
\item Construct $M_G$ column by column.  Prove the basis evaluation
identity $e(v)M_G(u)=G(u,v)e(v)$ from the monic-remainder identity,
under the two numerical polynomial-evaluation hypotheses.
\item Formalize the weight uniqueness lemma, the nonnegative exponent
checks for $F_t,G_t$, and the diagonal change-of-basis identity.  Negative
powers can be encoded as a finite Laurent coefficient table with a signed
starting exponent.  Alternatively multiply the identity by a computed
power of $t$ and work exclusively with ordinary polynomial coefficients.
\item Prove the explicit cyclic-matrix determinant, obtaining the nonzero
constant term.  Compare coefficients in the rescaling identity to derive
the norm degree and leading coefficient.
\item Prove the Vandermonde determinant and the fiber-divisibility lemma
using the explicit matrix $U$.  Prove the successive factor-removal
lemma by the triangular coefficient argument.  Apply it to the sorted
input point list to obtain the zero-count inequality.
\item Prove the finite count of monomials and the residue-permutation
identity.  Define a matrix presentation of a proper subspace, implement
rational row elimination with at most the matrix dimensions many pivot
steps, and build the annihilating row.  This gives the subspace and
finite-cover corollaries.
\end{enumerate}

For the object-language proof, every loop invariant can be expressed as
a Boolean equation about a bounded prefix: coefficient identities,
nonzero pivots, the residual factorization, and the current degree bound.
Parameters such as a point, a coefficient table or an input list remain
free during the induction.  Universal quantification over those
parameters is introduced afterwards.  No induction over an unbounded
statement about all points or all polynomials is required.

The accompanying \path{WeightedNormTargets.lean} defines exact rational
tables and states the propositions \path{NormDegreeGoal},
\path{FiberPowerGoal}, \path{CommonZerosGoal},
\path{WeightedBasisCountGoal}, and \path{EventualDimensionGoal}.
The proper-subspace corollary additionally uses the finite row-elimination
construction specified above. An ambient Lean proof of these
statements verifies their mathematical algebra.  To establish their
asserted iPRA derivability, the coefficient and loop proofs still have
to be translated into the repository's embedded proof judgment, and
then transported through its checked theory translation.  This section
specifies that route and gives the human proof; it does not assert that
those embedded derivation objects have already been constructed.
